% Anhang: Referenz des FourierGrating-Plugins
% Quellen: python/plugins/fouriergrating/fg_plugin.py (Parameter, Aktionen), fg_config.py,
% fg_scan.py, fg_nearfield.py, fg_estimate.py, fg_optimize.py, fg_cad.py, plugin.json;
% Blocktypen aus dem Paket rtt-fusion (rtt_fusion/graph/grating.py, stack.py)
\chapter{Referenz: FourierGrating-Plugin}\label{kap:plugin-fouriergrating}
\index{FourierGrating}

Dieser Anhang listet alle Aktionen, Parameter, Ergebnisse und Dateien des Plugins
\plugin{FourierGrating}. Die Bedienung erklärt \autoref{kap:fouriergrating}.

\section{Steckbrief}\label{sec:fg-ref-steckbrief}

\begin{referenz}{>{\raggedright\arraybackslash}p{3.4cm}>{\raggedright\arraybackslash}X}
\toprule
Eigenschaft & Wert \\
\midrule
\endhead
Name & \plugin{FourierGrating} \\
Kennung (\param{id}) & \code{fg} \\
Version & 0.1.0 \\
Plugin-API & 1 \\
Einstiegspunkt & \code{fg\_plugin:FourierGratingPlugin} \\
Python-Pakete (\param{requires}) & \code{numpy}, \code{scipy}, \code{fouriergrating>=0.18} \\
Ordner & \datei{python/plugins/fouriergrating} (wird mit der Anwendung installiert) \\
Blocktypen & \block{fg.grating}, \block{fg.stack} (mit \code{rtt-fusion}) \\
\bottomrule
\end{referenz}

\section{Python-Anforderungen}\label{sec:fg-ref-python}

\begin{referenz}{>{\raggedright\arraybackslash}p{4.2cm}>{\raggedright\arraybackslash}X}
\toprule
Paket / Merkmal & Wofür \\
\midrule
\endhead
\code{numpy} & alle Aktionen (Scan-Kurven, Nahfeld, CSV) \\
\code{scipy} & \gui{Optimise} (L-BFGS-B aus \code{scipy.optimize}) \\
\code{fouriergrating} $\geq$ 0.18 & \gui{Solve}, \gui{Scan}, \gui{Near field}, \gui{Optimise},
  \gui{Check CAD layers}; ohne das Paket meldet jede dieser Aktionen, dass es fehlt \\
CAD-Slicer in \code{fouriergrating} (Build mit \code{FG\_ENABLE\_NURBFORGE}) & \gui{Check CAD
  layers} und die Prüfung in \gui{Layer from CAD bodies} \\
\code{rtt\_fusion.\allowbreak graph.\allowbreak remote} (Paket \code{rtt-fusion} mit Fernblöcken) & Bereitstellen und
  Auswerten der Blöcke \block{fg.grating} und \block{fg.stack}; das Modul
  \code{rtt\_fusion.graph.stack} für \block{fg.stack} \\
\bottomrule
\end{referenz}

Ohne \code{fouriergrating} funktionieren: \gui{Load configuration}, \gui{Save configuration},
\gui{Estimate}, \gui{Set Littrow angle}, \gui{Layer from CAD bodies} (ohne Prüfung) sowie die beiden
Blockaktionen. Ein Gitterblock mit gespeicherter Jones-Tabelle rechnet ebenfalls ohne
\code{fouriergrating}.

% ---------------------------------------------------------------------------------------------
\section{Aktionen}\label{sec:fg-ref-aktionen}

Die Aktionen stehen im Menü \menu{Plugins,FourierGrating}, als Knöpfe im Panel \gui{Plugin} und in der
Werkzeugspalte -- in dieser Reihenfolge (alphabetisch nach dem internen Namen). Die beiden
Blockaktionen stehen nur im Kontextmenü eines Blocks \block{fg.grating} oder \block{fg.stack} im
Node-Editor.

\begin{referenz}{>{\raggedright\arraybackslash}p{3.6cm}>{\raggedright\arraybackslash}X>{\raggedright\arraybackslash}p{4.2cm}}
\toprule
Aktion & Wirkung & Liest / schreibt \\
\midrule
\endhead
\gui{Check CAD layers} & Geometrieprüfung und Bandzerlegung jeder \code{cad}-Schicht durch FourierGrating
  & Tabellen \enquote{Bands}, \enquote{Findings}; \gui{Preview SVG} \\
\gui{Estimate} & Zahl der Rechnungen, Matrixgröße, grober Speicherbedarf (Faustregel)
  & Statusleiste; Warnung im Protokoll \\
\gui{Layer from CAD bodies} & schreibt die gewählten Körper als STEP und hängt oben eine
  \code{cad}-Schicht an; prüft sie, wenn möglich & \gui{Bodies}, \gui{STEP file}; Parameter
  \param{layers} \\
\gui{Load configuration} & liest eine Konfiguration (.json) ins Panel & \gui{Configuration} \\
\gui{Near field} & Feld am ersten Punkt der Beleuchtung: vertikaler und horizontaler Schnitt &
  Bilder / Kurve; \gui{Also write .npz} \\
\gui{Optimise} & L-BFGS-B mit exakten Gradienten: Parameter innerhalb ihrer Grenzen für die
  Zielfunktion, am ersten Punkt & Kurve \enquote{objective L}, Tabelle \enquote{Parameters};
  Panel (mit \gui{Apply the optimum to the panel}) \\
\gui{Save configuration} & schreibt die Konfiguration des Panels als .json & \gui{Save as}, sonst
  \gui{Configuration} \\
\gui{Scan} & alle Punkte der Listen und Bereiche von $\lambda$, Winkeln, Polarisation und eines
  Sweeps: Effizienz je Ordnung als Kurven & Kurven; \gui{Also write CSV} (\gui{Scan}) \\
\gui{Set Littrow angle} & $\varphi$ für die \gui{Littrow order} (reflektiert, zurück in die
  Einfallsrichtung) bei der ersten Wellenlänge und dem ersten $\psi$ & Parameter \param{phi\_deg} \\
\gui{Solve} & Effizienzen jeder Ordnung am ersten Punkt der Beleuchtung & Tabelle \enquote{Orders},
  Balken; \gui{Solve: also write CSV} \\
\midrule
\gui{Store panel configuration in block} (Blockaktion) & die Konfiguration des Panels wird die des
  Blocks; Dateiverweis geleert, Jones-Tabelle eines Gitterblocks verworfen & Blockparameter
  \param{config}, \param{fouriergrating}, \param{table}, \param{table\_data}; Felder des
  Blockdialogs \\
\gui{Load block configuration into panel} (Blockaktion) & zeigt die Konfiguration des Blocks im
  Panel & Panel; \gui{Configuration} wird geleert \\
\bottomrule
\end{referenz}

% ---------------------------------------------------------------------------------------------
\section{Parameter}\label{sec:fg-ref-parameter}

Die Tabellen geben für jeden Parameter den Schlüssel (so steht er im Dokument und in
\code{ctx.params}), die Beschriftung im Panel, den Typ, den Standardwert, Einheit und Wertebereich
sowie die Bedeutung an. \enquote{frei} bedeutet: das Plugin legt keine Grenze fest. Die Gruppen
stehen in der Reihenfolge, in der das Panel sie zeigt.

Typen: \emph{Zahl} (Drehfeld mit Nachkommastellen), \emph{Ganzzahl}, \emph{Auswahl} (Klappliste),
\emph{Text} (Eingabezeile), \emph{Schalter}, \emph{Datei} (Pfad mit Dateidialog; \enquote{speichern}:
Dialog zum Speichern), \emph{Körper} (\gui{Use selection}), \emph{Liste} (Einträge mit \gui{Add},
\gui{Edit…}, \gui{Remove}, $\uparrow$, $\downarrow$), \emph{Tabelle} (mit \gui{Add row}, \gui{Remove
row}), \emph{JSON} (Textfeld, ungültiges JSON rot).

\subsection{Ohne Gruppe: Materialien und Schichten}

\begin{referenz}{>{\raggedright\arraybackslash}p{3.1cm}>{\raggedright\arraybackslash}p{2.7cm}>{\raggedright\arraybackslash}p{2.3cm}>{\raggedright\arraybackslash}X}
\toprule
Schlüssel / Beschriftung & Typ / Standard & Einheit / Bereich & Bedeutung \\
\midrule
\endhead
\param{materials}\newline\gui{Materials} & Liste (Einträge siehe unten); Standard: \code{air}
  ($n = 1$), \code{glass} ($n = 1{,}5$) & -- & alle Materialien; ein Eintrag wird in der Liste mit
  seinem \gui{Name} angezeigt \\
\param{layers}\newline\gui{Layers (bottom to top)} & Liste (Einträge siehe unten); Standard: leer &
  -- & Schichtstapel von unten (Substrat) nach oben (Superstrat); ein Eintrag wird mit seinem
  \gui{Type} angezeigt \\
\bottomrule
\end{referenz}

\paragraph{Felder eines Materials} (Dialog \gui{Edit entry} der Liste \gui{Materials}):

\begin{referenz}{>{\raggedright\arraybackslash}p{3.1cm}>{\raggedright\arraybackslash}p{2.7cm}>{\raggedright\arraybackslash}p{2.3cm}>{\raggedright\arraybackslash}X}
\toprule
Schlüssel / Beschriftung & Typ / Standard & Einheit / Bereich & Bedeutung \\
\midrule
\endhead
\param{name}\newline\gui{Name} & Text; \code{air} & -- & Name, unter dem Superstrat, Substrat und
  Schichten das Material nennen \\
\param{kind}\newline\gui{Given as} & Auswahl; \code{n} & \code{n}, \code{json} & \code{n}: Brechzahl
  $n + \mathrm{i}k$; \code{json}: jede andere Form (Permittivität, Tensoren, dispersive Modelle) \\
\param{n}\newline\gui{n} & Zahl (6 Stellen); 1.0 & frei & Realteil der Brechzahl (bei \code{n}) \\
\param{k}\newline\gui{k} & Zahl (6 Stellen); 0.0 & frei & Imaginärteil (Absorption); 0: das Material
  wird als reelle Zahl geschrieben \\
\param{json}\newline\gui{Material as JSON (json)} & JSON; leer (\code{null}) & -- & das Material in
  der Form der Konfigurationsdatei (bei \code{json}) \\
\bottomrule
\end{referenz}

\paragraph{Felder einer Schicht} (Dialog \gui{Edit entry} der Liste \gui{Layers (bottom to top)}):

\begin{referenz}{>{\raggedright\arraybackslash}p{3.1cm}>{\raggedright\arraybackslash}p{2.7cm}>{\raggedright\arraybackslash}p{2.3cm}>{\raggedright\arraybackslash}X}
\toprule
Schlüssel / Beschriftung & Typ / Standard & Einheit / Bereich & Bedeutung \\
\midrule
\endhead
\param{type}\newline\gui{Type} & Auswahl; \code{homogeneous} & \code{homogeneous},
  \code{lamellar1d}, \code{json} & Art der Schicht; \code{json}: jeder andere Schichttyp von
  FourierGrating, als JSON bearbeitet \\
\param{thickness}\newline\gui{Thickness} & Zahl (6 Stellen); 0.1 & \si{\micro\metre}; $\geq 0$ & Dicke
  (homogen und Lamellen) \\
\param{material}\newline\gui{Material (homogeneous)} & Text; \code{air} & -- & Material einer
  homogenen Schicht \\
\param{lamellae}\newline\gui{Lamellae (lamellar1d)} & Tabelle; leer & -- & Lamellen: jede von
  \gui{from x} bis zum \gui{from x} der nächsten Zeile, die letzte bis \gui{End} \\
\param{end}\newline\gui{End (lamellar1d)} & Zahl (6 Stellen); 0.0 & \si{\micro\metre}; frei & Ende der
  letzten Lamelle; muss gleich \gui{Period x} sein (nicht geprüft) \\
\param{json}\newline\gui{Layer as JSON (type json)} & JSON; leer & -- & die Schicht in der Form der
  Konfigurationsdatei (bei \code{json}) \\
\bottomrule
\end{referenz}

\paragraph{Spalten der Tabelle \gui{Lamellae (lamellar1d)}:}

\begin{referenz}{>{\raggedright\arraybackslash}p{3.1cm}>{\raggedright\arraybackslash}p{2.7cm}>{\raggedright\arraybackslash}p{2.3cm}>{\raggedright\arraybackslash}X}
\toprule
Schlüssel / Kopf & Typ / Standard & Einheit / Bereich & Bedeutung \\
\midrule
\endhead
\param{from}\newline\gui{from x} & Zahl (6 Stellen); 0.0 & \si{\micro\metre}; frei & Beginn der
  Lamelle; die erste Zeile beginnt bei 0 \\
\param{material}\newline\gui{material} & Text; \code{air} & -- & Material der Lamelle \\
\bottomrule
\end{referenz}

\subsection{Gruppe File}

\begin{referenz}{>{\raggedright\arraybackslash}p{3.1cm}>{\raggedright\arraybackslash}p{2.7cm}>{\raggedright\arraybackslash}p{2.3cm}>{\raggedright\arraybackslash}X}
\toprule
Schlüssel / Beschriftung & Typ / Standard & Einheit / Bereich & Bedeutung \\
\midrule
\endhead
\param{config\_file}\newline\gui{Configuration} & Datei; leer & Filter \enquote{FourierGrating
  configuration (*.json)} & Datei, die \gui{Load configuration} ins Panel liest; Ziel von
  \gui{Save configuration}, wenn \gui{Save as} leer ist \\
\param{save\_as}\newline\gui{Save as} & Datei (speichern); leer & wie oben & Ziel von \gui{Save
  configuration} \\
\bottomrule
\end{referenz}

\subsection{Gruppe Structure}

\begin{referenz}{>{\raggedright\arraybackslash}p{3.1cm}>{\raggedright\arraybackslash}p{2.7cm}>{\raggedright\arraybackslash}p{2.3cm}>{\raggedright\arraybackslash}X}
\toprule
Schlüssel / Beschriftung & Typ / Standard & Einheit / Bereich & Bedeutung \\
\midrule
\endhead
\param{dimension}\newline\gui{Dimension} & Auswahl; \code{1D} & \code{1D}, \code{2D} & Liniengitter
  (periodisch in x) oder Kreuzgitter (x und z) \\
\param{solver}\newline\gui{Solver} & Auswahl; \code{rcwa} & \code{rcwa}, \code{rk} & \code{rcwa}:
  Scheiben (Fourier-Modal-Methode); \code{rk}: differentielle Methode (implizites Runge-Kutta) \\
\param{nx}\newline\gui{Orders Nx} & Ganzzahl; 10 & 0 bis 200 & Fourier-Ordnungen $-N_x \dots N_x$ \\
\param{nz}\newline\gui{Orders Nz} & Ganzzahl; 0 & 0 bis 200 & Fourier-Ordnungen $-N_z \dots N_z$
  (2D) \\
\param{period\_x}\newline\gui{Period x} & Zahl (6 Stellen); 1.0 & \si{\micro\metre}; $\geq 0$ &
  Periode $\Lambda_x$ \\
\param{period\_z}\newline\gui{Period z} & Zahl (6 Stellen); 0.0 & \si{\micro\metre}; $\geq 0$ &
  Periode $\Lambda_z$ (2D: $> 0$ nötig) \\
\param{superstrate}\newline\gui{Superstrate} & Text; \code{air} & Materialname & Medium oben, aus
  dem das Licht einfällt (verlustfrei) \\
\param{substrate}\newline\gui{Substrate} & Text; \code{glass} & Materialname & Medium unten \\
\param{solve\_csv}\newline\gui{Solve: also write CSV} & Datei (speichern); leer & Filter
  \enquote{CSV (*.csv)} & \gui{Solve} schreibt die Ordnungstabelle zusätzlich hierhin \\
\bottomrule
\end{referenz}

\subsection{Gruppe Incidence}

\begin{referenz}{>{\raggedright\arraybackslash}p{3.1cm}>{\raggedright\arraybackslash}p{2.7cm}>{\raggedright\arraybackslash}p{2.3cm}>{\raggedright\arraybackslash}X}
\toprule
Schlüssel / Beschriftung & Typ / Standard & Einheit / Bereich & Bedeutung \\
\midrule
\endhead
\param{lambda\_um}\newline\gui{$\lambda$ / \si{\micro\metre}} & Text; \code{0.6328} &
  \si{\micro\metre}; Zahl, Liste oder Bereich & Vakuumwellenlänge: \code{0.6}, \code{0.5, 0.6} oder
  \code{0.5:0.7:21} (von:bis:Schritte) \\
\param{phi\_deg}\newline\gui{$\varphi$ / deg} & Text; \code{0} & Grad; wie oben & Polarwinkel gegen
  die Normale (y-Achse) \\
\param{psi\_deg}\newline\gui{$\psi$ / deg} & Text; \code{0} & Grad; wie oben & Azimut der
  Einfallsebene gegen x \\
\param{delta\_deg}\newline\gui{$\delta$ / deg} & Text; \code{0} & Grad; wie oben & Polarisation:
  0 = s (TE), 90 = p (TM) \\
\param{littrow\_order}\newline\gui{Littrow order} & Ganzzahl; $-1$ & $-50$ bis 50 & Ordnung, für die
  \gui{Set Littrow angle} $\varphi$ setzt \\
\bottomrule
\end{referenz}

Ein leeres Beleuchtungsfeld lässt den Schlüssel in der Konfiguration weg.

\subsection{Gruppe Scan}

\begin{referenz}{>{\raggedright\arraybackslash}p{3.1cm}>{\raggedright\arraybackslash}p{2.7cm}>{\raggedright\arraybackslash}p{2.3cm}>{\raggedright\arraybackslash}X}
\toprule
Schlüssel / Beschriftung & Typ / Standard & Einheit / Bereich & Bedeutung \\
\midrule
\endhead
\param{scan\_x}\newline\gui{x axis} & Auswahl; \code{auto} & \code{auto}, \code{lambda\_um},
  \code{phi\_deg}, \code{psi\_deg}, \code{delta\_deg}, \code{sweep\_value} & Größe auf der x-Achse;
  \code{auto}: die mit den meisten Werten (bei Gleichstand $\lambda$); die übrigen variierenden
  Größen ergeben eigene Kurven \\
\param{scan\_min}\newline\gui{Orders from efficiency} & Zahl (4 Stellen); 0.01 & 0 bis 1 & eine
  Ordnung bekommt eine Kurve, wenn sie diese Effizienz irgendwo erreicht \\
\param{scan\_max\_curves}\newline\gui{At most curves} & Ganzzahl; 24 & 1 bis 500 & höchstens so
  viele Kurven; weitere werden mit Warnung weggelassen \\
\param{scan\_csv}\newline\gui{Also write CSV} & Datei (speichern); leer & Filter \enquote{CSV
  (*.csv)} & jede Ordnung an jedem Punkt, wie \datei{efficiencies.csv} von FourierGrating \\
\bottomrule
\end{referenz}

\subsection{Gruppe Near field}

\begin{referenz}{>{\raggedright\arraybackslash}p{3.1cm}>{\raggedright\arraybackslash}p{2.7cm}>{\raggedright\arraybackslash}p{2.3cm}>{\raggedright\arraybackslash}X}
\toprule
Schlüssel / Beschriftung & Typ / Standard & Einheit / Bereich & Bedeutung \\
\midrule
\endhead
\param{nf\_show}\newline\gui{Show} & Auswahl; \code{|E|$^2$} & siehe \autoref{sec:fg-ref-nahfeld} &
  dargestellte Größe \\
\param{nf\_y\_from}\newline\gui{y from} & Zahl (6 Stellen); $-0.2$ & \si{\micro\metre}; frei & unteres
  Ende des vertikalen Schnitts (0 = Unterkante des Stapels) \\
\param{nf\_y\_to}\newline\gui{y to} & Zahl (6 Stellen); 0.5 & \si{\micro\metre}; frei & oberes Ende
  des vertikalen Schnitts \\
\param{nf\_horizontal\_y}\newline\gui{Horizontal cut at y} & Zahl (6 Stellen); 0.1 &
  \si{\micro\metre}; frei & Höhe des horizontalen Schnitts \\
\param{nf\_vertical\_z}\newline\gui{Vertical cut at z (2D)} & Zahl (6 Stellen); 0.0 &
  \si{\micro\metre}; frei & Lage des vertikalen Schnitts in z (2D) \\
\param{nf\_nx}\newline\gui{Points x} & Ganzzahl; 400 & 2 bis 4000 & Punkte in x \\
\param{nf\_ny}\newline\gui{Points y} & Ganzzahl; 300 & 2 bis 4000 & Punkte in y \\
\param{nf\_nz}\newline\gui{Points z (2D)} & Ganzzahl; 200 & 2 bis 4000 & Punkte in z (2D) \\
\param{nf\_periods}\newline\gui{Periods} & Ganzzahl; 2 & 1 bis 20 & gezeigte Perioden \\
\param{nf\_npz}\newline\gui{Also write .npz} & Datei (speichern); leer & Filter \enquote{NumPy
  archive (*.npz)} & jede Größe beider Schnitte mit ihren Achsen \\
\bottomrule
\end{referenz}

{\sloppy Die Werte \param{nf\_horizontal\_y}, \param{nf\_vertical\_z}, \param{nf\_nx}, \param{nf\_ny},
\param{nf\_nz}, \param{nf\_periods} und der Bereich \param{nf\_y\_from} bis \param{nf\_y\_to} entsprechen den
Schlüsseln \code{horizontal\_y\_um}, \code{vertical\_z\_um}, \code{nx}, \code{ny}, \code{nz},
\code{periods} und \code{y\_range\_um} des Blocks \code{nearfield}. Sie werden nur geschrieben, wenn die
geladene Datei sie hatte oder sie vom Standard abweichen; \gui{Near field} rechnet immer mit allen.\par}

\subsection{Gruppe Optimisation}

\begin{referenz}{>{\raggedright\arraybackslash}p{3.1cm}>{\raggedright\arraybackslash}p{2.7cm}>{\raggedright\arraybackslash}p{2.3cm}>{\raggedright\arraybackslash}X}
\toprule
Schlüssel / Beschriftung & Typ / Standard & Einheit / Bereich & Bedeutung \\
\midrule
\endhead
\param{opt\_goal}\newline\gui{Goal} & Auswahl; \code{maximise} & \code{maximise}, \code{minimise} &
  Richtung der Optimierung \\
\param{opt\_terms}\newline\gui{Objective L = $\Sigma$ weight $\cdot$ efficiency} & Tabelle; leer &
  Spalten siehe unten & Glieder der Zielfunktion $L = \sum w\,\eta(\text{side}, n, m)$ \\
\param{opt\_params}\newline\gui{Parameters} & Tabelle; leer & Spalten siehe unten & Zahlen der
  Konfiguration per JSON-Zeiger, mit Grenzen \\
\param{opt\_iterations}\newline\gui{Iterations} & Ganzzahl; 50 & 1 bis 10000 & höchstens so viele
  Iterationen von L-BFGS-B \\
\param{opt\_apply}\newline\gui{Apply the optimum to the panel} & Schalter; an & -- & schreibt das
  Optimum nach dem Ende ins Panel \\
\bottomrule
\end{referenz}

\paragraph{Spalten von \gui{Objective L = $\Sigma$ weight $\cdot$ efficiency}:}

\begin{referenz}{>{\raggedright\arraybackslash}p{3.1cm}>{\raggedright\arraybackslash}p{2.7cm}>{\raggedright\arraybackslash}p{2.3cm}>{\raggedright\arraybackslash}X}
\toprule
Schlüssel / Kopf & Typ / Standard & Einheit / Bereich & Bedeutung \\
\midrule
\endhead
\param{side}\newline\gui{side} & Auswahl; \code{R} & \code{R}, \code{T} & reflektierte oder
  transmittierte Ordnung \\
\param{n}\newline\gui{n} & Ganzzahl; $-1$ & $-1000$ bis 1000 & Ordnung in x \\
\param{m}\newline\gui{m} & Ganzzahl; 0 & $-1000$ bis 1000 & Ordnung in z (1D: 0) \\
\param{weight}\newline\gui{weight} & Zahl (4 Stellen); 1.0 & frei & Gewicht $w$; Zeilen derselben
  Ordnung werden addiert \\
\bottomrule
\end{referenz}

\paragraph{Spalten von \gui{Parameters}:}

\begin{referenz}{>{\raggedright\arraybackslash}p{3.1cm}>{\raggedright\arraybackslash}p{2.7cm}>{\raggedright\arraybackslash}p{2.3cm}>{\raggedright\arraybackslash}X}
\toprule
Schlüssel / Kopf & Typ / Standard & Einheit / Bereich & Bedeutung \\
\midrule
\endhead
\param{pointer}\newline\gui{JSON pointer} & Text; \code{/layers/0/\allowbreak thickness\_um} & RFC 6901 & Zahl
  der Konfiguration, z.\,B. \code{/layers/0/breakpoints\_um/1}, \code{/period\_um/x} \\
\param{lower}\newline\gui{lower} & Zahl (6 Stellen); 0.0 & frei & Untergrenze \\
\param{upper}\newline\gui{upper} & Zahl (6 Stellen); 1.0 & frei & Obergrenze; der Startwert muss
  zwischen beiden liegen \\
\bottomrule
\end{referenz}

\subsection{Gruppe CAD layer}

\begin{referenz}{>{\raggedright\arraybackslash}p{3.1cm}>{\raggedright\arraybackslash}p{2.7cm}>{\raggedright\arraybackslash}p{2.3cm}>{\raggedright\arraybackslash}X}
\toprule
Schlüssel / Beschriftung & Typ / Standard & Einheit / Bereich & Bedeutung \\
\midrule
\endhead
\param{cad\_bodies}\newline\gui{Bodies} & Körper (mehrere); leer & Art \code{shell} & geschlossene
  Körper einer Elementarzelle; Material aus dem Attribut \code{material} oder dem Namen; CAD $+Z$ ist
  die Stapelrichtung \\
\param{cad\_unit}\newline\gui{Bodies drawn in} & Auswahl; \si{\micro\metre} & \si{\micro\metre},
  \code{mm} & Einheit der gezeichneten Körper \\
\param{cad\_step}\newline\gui{STEP file} & Datei (speichern); leer & Filter \enquote{STEP (*.step
  *.stp)} & wohin die Körper geschrieben werden; die Schicht verweist darauf \\
\param{cad\_background}\newline\gui{Around the bodies} & Auswahl; \code{superstrate} &
  \code{superstrate}, \code{substrate} & Material zwischen den Körpern \\
\param{cad\_slices}\newline\gui{RCWA slices per band} & Ganzzahl; 30 & 1 bis 1000 & Scheiben je Band
  der RCWA \\
\param{cad\_preview\_y}\newline\gui{Preview at height} & Zahl (6 Stellen); 0.0 & \si{\micro\metre};
  frei & Höhe des Vorschauschnitts \\
\param{cad\_preview}\newline\gui{Preview SVG} & Datei (speichern); leer & Filter \enquote{SVG
  (*.svg)} & \gui{Check CAD layers} schreibt den Schnitt der ersten \code{cad}-Schicht in dieser
  Höhe \\
\bottomrule
\end{referenz}

\subsection{Gruppe Advanced und verborgene Parameter}

\begin{referenz}{>{\raggedright\arraybackslash}p{3.1cm}>{\raggedright\arraybackslash}p{2.7cm}>{\raggedright\arraybackslash}p{2.3cm}>{\raggedright\arraybackslash}X}
\toprule
Schlüssel / Beschriftung & Typ / Standard & Einheit / Bereich & Bedeutung \\
\midrule
\endhead
\param{advanced}\newline\gui{Further settings (JSON)} & JSON; \code{\{\}} & -- & alle übrigen
  Schlüssel der Konfiguration (Solver-Optionen, \code{sweep}, \code{output}, weitere Schlüssel von
  \code{incidence}, \code{truncation}, \code{nearfield}), wie in der Datei \\
\param{\_order} & JSON, verborgen; \code{[]} & -- & Reihenfolge der Schlüssel der geladenen Datei
  (wird beim Speichern eingehalten) \\
\param{\_present} & JSON, verborgen; \code{null} & -- & welche Schlüssel die geladene Datei hatte
  (nur diese oder geänderte werden geschrieben) \\
\bottomrule
\end{referenz}

% ---------------------------------------------------------------------------------------------
\section{Panel und Konfigurationsdatei}\label{sec:fg-ref-zuordnung}

So entsprechen die Felder des Panels den Schlüsseln der Konfiguration von FourierGrating:

\begin{referenz}{>{\raggedright\arraybackslash}p{5.0cm}>{\raggedright\arraybackslash}X}
\toprule
Panel & Konfiguration \\
\midrule
\endhead
\param{dimension}, \param{solver} & \code{dimension}, \code{solver} \\
\param{nx}, \param{nz} & \code{truncation.Nx}, \code{truncation.Nz} \\
\param{period\_x}, \param{period\_z} & \code{period\_um.x}, \code{period\_um.z} \\
\param{superstrate}, \param{substrate} & \code{superstrate}, \code{substrate} \\
\param{lambda\_um}, \param{phi\_deg}, \param{psi\_deg}, \param{delta\_deg} & \code{incidence.lambda\_um}
  usw.; Liste \code{[...]}, Bereich \code{\{\textquotedbl{}from\textquotedbl{}, \textquotedbl{}to\textquotedbl{}, \textquotedbl{}steps\textquotedbl{}\}} \\
Material mit \gui{Given as} = \code{n} & \code{materials.<Name>.n}: Zahl (bei $k = 0$) oder
  \code{[n, k]} \\
Material mit \gui{Given as} = \code{json} & \code{materials.<Name>}: der JSON-Inhalt \\
Schicht \code{homogeneous} & \code{\{\textquotedbl{}type\textquotedbl{}: \textquotedbl{}homogeneous\textquotedbl{}, \textquotedbl{}thickness\_um\textquotedbl{}, \textquotedbl{}material\textquotedbl{}\}} \\
Schicht \code{lamellar1d} & \code{\{\textquotedbl{}type\textquotedbl{}: \textquotedbl{}lamellar1d\textquotedbl{}, \textquotedbl{}thickness\_um\textquotedbl{}, \textquotedbl{}breakpoints\_um\textquotedbl{},
  \textquotedbl{}materials\textquotedbl{}\}}; \code{breakpoints\_um} = alle \gui{from x} und \gui{End} \\
Schicht \code{json} & der JSON-Inhalt \\
Gruppe \gui{Near field} & \code{nearfield} (siehe oben) \\
\param{advanced} & alle übrigen Schlüssel, unverändert \\
\bottomrule
\end{referenz}

Beim Laden erscheinen homogene und Lamellenschichten nur dann als Felder, wenn sie genau die oben
genannten Schlüssel haben; Materialien nur, wenn sie allein aus \code{n} (Zahl oder Paar) bestehen.
Alles andere erscheint als JSON-Eintrag.

% ---------------------------------------------------------------------------------------------
\section{Ergebnisse}\label{sec:fg-ref-ergebnisse}

\begin{referenz}{>{\raggedright\arraybackslash}p{2.8cm}>{\raggedright\arraybackslash}p{4.2cm}>{\raggedright\arraybackslash}X}
\toprule
Aktion & Ergebnissatz & Inhalt \\
\midrule
\endhead
\gui{Solve} & \enquote{FourierGrating $\lambda$ = … \si{\micro\metre}} & Tabelle \enquote{Orders}
  (Spalten \code{side}, \code{n}, \code{m}, \code{efficiency}, \code{s part}, \code{p part});
  Balken \enquote{Efficiency per order} (y-Achse \code{efficiency}); Protokoll: $\Sigma R$, $\Sigma T$,
  $A$, Energiebilanz, Warnungen \\
\gui{Scan} & \enquote{FourierGrating scan over …} (Achsenbeschriftung) & Kurven $\Sigma$R, $\Sigma$T und
  je Ordnung (\code{R n} bzw. \code{R n,m}, mit Zusatz der übrigen variierenden Größen); y-Achse
  \code{efficiency} \\
\gui{Near field} & \enquote{FourierGrating near field $\lambda$ = … \si{\micro\metre}} & Bild bzw.
  Kurve \enquote{<Größe> – vertical cut} und \enquote{<Größe> – horizontal cut}; Beschreibung: die
  Normierung \\
\gui{Optimise} & \enquote{FourierGrating optimisation} & Kurve \enquote{objective L} (x:
  \code{evaluation}, y: \code{L}); Tabelle \enquote{Parameters} (\code{parameter}, \code{start},
  \code{optimum}, \code{lower}, \code{upper}) \\
\gui{Check CAD layers}, \gui{Layer from CAD bodies} & \enquote{FourierGrating CAD layers} & Tabelle
  \enquote{Bands} (\code{layer}, \code{y from / }\si{\micro\metre}, \code{y to /
  }\si{\micro\metre}, \code{kind}, \code{area fractions}); Tabelle \enquote{Findings}
  (\code{layer}, \code{severity}, \code{code}, \code{message}, \code{where / }\si{\micro\metre}) \\
\bottomrule
\end{referenz}

\subsection{Größen des Nahfelds}\label{sec:fg-ref-nahfeld}

\begin{referenz}{>{\raggedright\arraybackslash}p{3.6cm}>{\raggedright\arraybackslash}p{2.6cm}>{\raggedright\arraybackslash}p{2.4cm}>{\raggedright\arraybackslash}X}
\toprule
\gui{Show} & Feld von FourierGrating & Farbskala & Darstellung \\
\midrule
\endhead
\code{|E|$^2$} & \code{absE2} & \code{inferno} & Wert \\
\code{|H|$^2$} & \code{absH2} & \code{inferno} & Wert \\
\code{Sy (power flow)} & \code{Sy} & \code{coolwarm} & Wert, Skala symmetrisch um 0 \\
\code{|Ex|}, \code{|Ey|}, \code{|Ez|}, \code{|Hx|}, \code{|Hy|}, \code{|Hz|} & \code{Ex} \dots{}
  \code{Hz} & \code{inferno} & Betrag \\
\code{Re Ex}, \code{Re Ey}, \code{Re Ez}, \code{Re Hx}, \code{Re Hy}, \code{Re Hz} & \code{Ex} \dots{}
  \code{Hz} & \code{coolwarm} & Realteil, Skala symmetrisch um 0 \\
\code{material} & \code{material} & \code{viridis} & Materialindex \\
\bottomrule
\end{referenz}

Normierung: $|E_\text{inc}| = 1$, das Magnetfeld als $Z_0 H$.

% ---------------------------------------------------------------------------------------------
\section{Exportierte Dateien}\label{sec:fg-ref-dateien}

\begin{description}
  \item[Konfiguration] (\gui{Save configuration}) JSON mit zwei Leerzeichen Einrückung, Schlüssel in
    der Reihenfolge der geladenen Datei; Format nach \datei{docs/config\_format.md} von
    FourierGrating.
  \item[Ordnungstabelle] (\gui{Solve: also write CSV}) eine Zeile je Ordnung, Spalten (Kopfzeile) \code{side}, \code{n}, \code{m}, \code{efficiency}, \code{eff\_s}, \code{eff\_p}; Zahlen mit voller Genauigkeit.
  \item[Scan] (\gui{Also write CSV}) alle Spalten des Scan-Ergebnisses von FourierGrating in seiner
    Reihenfolge -- wie \datei{efficiencies.csv} von FourierGrating (u.\,a. \code{lambda\_um},
    \code{phi\_deg}, \code{psi\_deg}, \code{delta\_deg}, \code{side}, \code{n}, \code{m},
    \code{efficiency}, \code{eff\_s}, \code{eff\_p}, bei einem Sweep dessen Wert); komplexe Spalten
    als \code{<Name>\_re} und \code{<Name>\_im}; eine Zeile je Ordnung und Punkt.
  \item[Nahfeld] (\gui{Also write .npz}) komprimiertes NumPy-Archiv: jedes Feld jedes Schnitts unter
    \code{<Schnitt>/<Name>} (komplexe Felder unverändert), dazu \code{<Schnitt>/meta} als JSON-Text;
    Schnitte \code{vertical} und \code{horizontal} (und die Angaben zur Beleuchtung).
  \item[STEP] (\gui{STEP file}) die gewählten Körper, je einer benannt nach seinem Objekt, in der
    gewählten Reihenfolge.
  \item[Vorschau] (\gui{Preview SVG}) der Schnitt der ersten \code{cad}-Schicht in der Höhe
    \gui{Preview at height}.
\end{description}

\subsection{Die angelegte CAD-Schicht}

\gui{Layer from CAD bodies} hängt diesen Eintrag (als Schicht vom Typ \code{json}) an:

\begin{lstlisting}
{ "type": "cad",
  "file": "<absoluter Pfad der STEP-Datei>",
  "unit": "um",                      (oder "mm")
  "axes": { "y": "+Z", "x": "+X" },
  "background": "superstrate",       (oder "substrate")
  "bodies": [ { "index": 0, "material": "si" }, ... ],
  "slicing": { "rcwa_slices": 30 } }
\end{lstlisting}

% ---------------------------------------------------------------------------------------------
\section{Formeln der Aufwandsschätzung}\label{sec:fg-ref-schaetzung}

\begin{referenz}{>{\raggedright\arraybackslash}p{4.0cm}>{\raggedright\arraybackslash}X}
\toprule
Größe & Berechnung \\
\midrule
\endhead
Scanpunkte & Produkt der Zahl der Werte von $\lambda$, $\varphi$, $\psi$, $\delta$ (Zahl: 1, Liste:
  Länge, Bereich: \code{steps}) \\
Rechnungen & Scanpunkte $\times$ Sweep-Werte \\
Fourier-Glieder $n_f$ & $2N_x + 1$ (1D), $(2N_x + 1)(2N_z + 1)$ (2D) \\
Matrixgröße & $n = 4 n_f$ \\
Speicher & $\min(\text{Scanpunkte}, \text{Kerne}) \times 12 \times n^2 \times 16$ Byte; bei
  \gui{Optimise} ein Punkt \\
Warnschwelle & 70\,\% des Arbeitsspeichers \\
\bottomrule
\end{referenz}

% ---------------------------------------------------------------------------------------------
\section{Workflow-Blöcke}\label{sec:fg-ref-bloecke}
\index{fg.grating}\index{fg.stack}

Die Blocktypen stammen aus dem Paket \code{rtt-fusion} und werden im Prozess des Plugins
ausgewertet. Die folgenden Angaben entsprechen der installierten Version des Pakets; andere Versionen
können weitere Parameter haben. Die Darstellungen sind R (Strahlen) und H (Felder). Allgemeines zu
Ports und Blöcken steht in \autoref{kap:workflow} und \autoref{kap:workflow-bloecke}.

\subsection{Gitter: fg.grating}

Periodische Probe (\enquote{Periodic sample (FourierGrating table): diffraction orders as channels}),
Kategorie Strahlteiler. Frühere Kennung \block{fusion.grating} (wird beim Laden umgesetzt).

\begin{referenz}{>{\raggedright\arraybackslash}p{2.6cm}>{\raggedright\arraybackslash}p{2.4cm}>{\raggedright\arraybackslash}X}
\toprule
Anschluss & Richtung, Art & Bedeutung \\
\midrule
\endhead
\code{in} & Eingang, Licht (R, H) & einfallendes Licht \\
\code{out} (\enquote{orders}) & Ausgang, Licht (R, H) & die weitergegebenen Ordnungen \\
Port \code{front} & \code{in}/\code{out} & Seite der Probe, an der Licht ein- und austritt \\
\bottomrule
\end{referenz}

\begin{referenz}{>{\raggedright\arraybackslash}p{3.0cm}>{\raggedright\arraybackslash}p{2.8cm}>{\raggedright\arraybackslash}X}
\toprule
Parameter & Typ / Standard & Bedeutung \\
\midrule
\endhead
\param{table} & Datei; leer & Jones-Tabelle (\code{raytatouille-result}-Datei); \gui{Check workflow}
  übernimmt sie in den Block (\param{table\_data}) und leert den Pfad \\
\param{fouriergrating} & Datei; leer & FourierGrating-Konfiguration; der Block baut seine Tabelle um
  die Strahlen (R) oder rechnet jede ebene Welle eines Felds (H); \gui{Check workflow} übernimmt sie
  in \param{config} \\
\param{config} & Daten; \code{\{\}} & die Konfiguration, die der Block hält (vor der Datei) \\
\param{table\_data} & Daten; \code{\{\}} & die Jones-Tabelle, die der Block hält (vor der Datei; hat
  Vorrang vor der Konfiguration) \\
\param{side} & Auswahl; \code{reflected} & \code{reflected} oder \code{transmitted}: die
  weitergegebenen Ordnungen (mit Konfiguration) \\
\param{orders} & Zahlenliste; leer & weitergegebene Ordnungen; leer: alle der Tabelle \\
\param{n\_in} & Zahl; 1.0 & Brechzahl des Einfallsmediums (R, gebaute Tabelle) \\
\param{sigma\_step} (\gui{Ray path}: \gui{$\sigma$ step of the table}) & Zahl; 0.01 ($\geq 10^{-5}$) &
  Abstand der Stützstellen der Tabelle in den Richtungskosinus $\sigma = n \sin\theta$ (kubisch
  interpoliert); kleiner ist genauer und braucht mehr Rechnungen \\
\param{sigma\_margin} (\gui{Ray path}: \gui{$\sigma$ margin of the table}) & Zahl; 0.01 ($\geq 0$) &
  wie weit die Tabelle über den $\sigma$-Bereich der Strahlen hinausreicht \\
\param{point} & Punkt; (0, 0, 0) \si{\milli\metre} & Lage der Probe \\
\param{normal} & Punkt; (0, 0, 1) & Flächennormale \\
\param{grating\_vector} & Punkt; (1, 0, 0) & lokale x-Richtung, senkrecht zu den Gitterlinien \\
\param{field\_transfer} (\gui{Field path}) & Auswahl; \code{exact} & \code{exact}: FourierGrating je
  ebener Welle; \code{interpolated}: aus einem geprüften Teilgitter des Spektrums, für breite
  Bündel viel schneller \\
\param{field\_threshold} (\gui{Field path}) & Zahl; $10^{-10}$ & ebene Wellen unter diesem Anteil
  der größten Amplitude entfallen (H) \\
\param{field\_grid\_max} (\gui{Field path}) & Ganzzahl; 1024 & höchstens so viele Gitterpunkte je
  Achse für ein Feld (H) \\
\bottomrule
\end{referenz}

Eine gebaute Tabelle umfasst die Wellenlänge der Strahlen $\pm$\,\qty{1.5}{\nano\metre} und ihren
$\sigma$-Bereich plus \param{sigma\_margin}, mindestens vier Stützstellen je Achse. Strahlen, deren
Ordnung evaneszent ist oder die außerhalb der Tabelle bzw. nahe einer Rayleigh-Anomalie liegen,
entfallen mit einer Protokollzeile. Alle Ordnungen behalten die Kohärenz des einfallenden Bündels.

Hält der Block eine Konfiguration, zeigt sein Dialog zusätzlich deren Zahlen:

\begin{referenz}{>{\raggedright\arraybackslash}p{2.4cm}>{\raggedright\arraybackslash}p{5.4cm}>{\raggedright\arraybackslash}X}
\toprule
Gruppe & Feld & Zeiger in die Konfiguration \\
\midrule
\endhead
\gui{Grating} & \gui{Period x} (2D auch \gui{Period z}), \si{\micro\metre} & \code{/config/period\_um/x} \\
\gui{Grating} & \gui{Fourier orders Nx ($-$N … N)} (2D auch Nz) & \code{/config/truncation/Nx} \\
\gui{Layers} & \gui{Layer $i$ (<Typ>): thickness}, \si{\micro\metre} & \code{/config/layers/<i-1>/thickness\_um} \\
\gui{Layers} & \gui{Layer $i$ (<Typ>): lamella edges}, \si{\micro\metre} & \code{/config/layers/<i-1>/breakpoints\_um} \\
\gui{Materials} & \gui{<Name>: n} bzw. \gui{<Name>: n, k} & \code{/config/materials/<Name>/n} \\
\bottomrule
\end{referenz}

\subsection{Schichtstapel: fg.stack}

\enquote{Layer stack (FourierGrating, exact per ray): coating, filter, mirror layers}. Licht, das auf
einer Seite eintritt, verlässt den Stapel reflektiert auf derselben und transmittiert auf der
anderen Seite. Licht von hinten sieht den Stapel umgedreht.

\begin{referenz}{>{\raggedright\arraybackslash}p{2.6cm}>{\raggedright\arraybackslash}p{2.4cm}>{\raggedright\arraybackslash}X}
\toprule
Anschluss & Richtung, Art & Bedeutung \\
\midrule
\endhead
\code{front\_in} & Eingang, Licht (R, H) & Licht auf die Vorderseite (Superstrat, in Richtung
  \param{normal}) \\
\code{front\_out} & Ausgang, Licht (R, H) & Licht, das die Vorderseite verlässt \\
\code{back\_in} & Eingang, Licht (R, H) & Licht auf die Rückseite (Substrat) \\
\code{back\_out} & Ausgang, Licht (R, H) & Licht, das die Rückseite verlässt \\
\bottomrule
\end{referenz}

\begin{referenz}{>{\raggedright\arraybackslash}p{3.0cm}>{\raggedright\arraybackslash}p{2.8cm}>{\raggedright\arraybackslash}X}
\toprule
Parameter & Typ / Standard & Bedeutung \\
\midrule
\endhead
\param{fouriergrating} & Datei; leer & Konfiguration mit nur homogenen Schichten, vom Substrat (hinten)
  nach oben zum Superstrat (vorn); \gui{Check workflow} übernimmt sie in \param{config} \\
\param{config} & Daten; \code{\{\}} & die Konfiguration, die der Block hält \\
\param{point} & Punkt; (0, 0, 0) \si{\milli\metre} & Lage der Ebene des Stapels \\
\param{normal} & Punkt; (0, 0, 1) & zur Vorderseite (Superstrat) \\
\param{field\_threshold} (\gui{Field path}) & Zahl; $10^{-10}$ & wie bei \block{fg.grating} \\
\param{field\_transfer} (\gui{Field path}) & Auswahl; \code{exact} & wie bei \block{fg.grating} \\
\param{field\_grid\_max} (\gui{Field path}) & Ganzzahl; 1024 & wie bei \block{fg.grating} \\
\bottomrule
\end{referenz}

Beleuchtung, Abbruch und Periode der Konfiguration setzt der Block selbst ($N = 0$, 1D). Der
Stapel ist ein dünnes Element: die seitliche Versetzung im Stapel (höchstens Dicke $\times
\tan\theta$) entfällt.
